% 2.3 局部 Langlands 对应的陈述（底本 p.21--22）
\subsection{局部 Langlands 对应的陈述}

到目前为止，我们还没有引入 Galois 表示与容许表示的 $L$-因子和 $\epsilon$-因子，
也没有引入（本质）平方可积、（本质）缓增与一般（generic）表示的概念。
不过我认为此刻已经可以陈述精确的局部 Langlands 对应（LLC）了。

\subsubsection{LLC 的陈述}

\begin{theorem}
存在唯一的一族集合间双射
\[
\mathrm{rec}_n\colon
\left\{\begin{array}{c}
\GL_n(K)\ \text{的不可约容许表示}\\
\text{（的等价类）}
\end{array}\right\}
\xrightarrow{\ \sim\ }
\left\{\begin{array}{c}
W_K\ \text{的}\ n\ \text{维 Frobenius 半单}\\
\text{Weil--Deligne 表示（的等价类）}
\end{array}\right\},
\]
满足下列诸条性质：
\begin{enumerate}[label=(\arabic*),leftmargin=2.2em]
\item 对 $K^{\times}$ 的每个不可约拟特征标 $\chi$，
\[
\mathrm{rec}_1(\chi)=\chi\circ\Art_K^{-1},
\]
其中 $\Art_K\colon K^{\times}\xrightarrow{\ \sim\ }W^{\mathrm{ab}}_K$。
\item 对 $\GL_{n_1}$、$\GL_{n_2}$ 的不可约容许表示 $\pi_1,\pi_2$，
\[
L(\pi_1\times\pi_2,s)=L\bigl(\mathrm{rec}_{n_1}(\pi_1)\otimes\mathrm{rec}_{n_2}(\pi_2),s\bigr),
\qquad
\epsilon(\pi_1\times\pi_2,s)=\epsilon\bigl(\mathrm{rec}_{n_1}(\pi_1)\otimes
\mathrm{rec}_{n_2}(\pi_2),s\bigr),
\]
其中 $L$-因子与 $\epsilon$-因子在第 3.6 小节中定义。
\item 对 $K^{\times}$ 的每个拟特征标 $\chi$ 与每个不可约 $\pi\in\Repadm_{\GL_n(K)}$，有
\[
\mathrm{rec}_n(\pi\chi)=\mathrm{rec}_n(\pi)\otimes\mathrm{rec}_1(\chi).
\]
\item 对每个中心特征标为 $\omega_{\pi}$ 的不可约 $\pi\in\Repadm_{\GL_n(K)}$，
\[
\det\circ\mathrm{rec}_n(\pi)=\mathrm{rec}_1(\omega_{\pi}).
\]
\item 对每个不可约 $\pi\in\Repadm_{\GL_n(K)}$，有
\[
\mathrm{rec}_n(\pi^{\vee})=\mathrm{rec}_n(\pi)^{\vee}.
\]
\end{enumerate}
\end{theorem}

\begin{remark}
局部 Langlands 对应 $\mathrm{rec}_n$ 保持两个重要的不变量：
\begin{itemize}[leftmargin=1.6em]
\item \emph{导子}：$c(\pi)=c(\mathrm{rec}_n(\pi))$，这来自与 $\epsilon$-因子的相容性。
\item \emph{深度}：$d(\pi)=d(\mathrm{rec}_n(\pi))$，由 Moy--Prasad 引入。
\end{itemize}
\end{remark}

\begin{remark}
由 Henniart 的工作，只需证明存在一族满足上述 (1)--(5) 的映射 $\mathrm{rec}_n$
（作用在超尖点表示上）。Henniart 的定理断言：一旦 $\mathrm{rec}_n$ 存在，
它必然是一族双射且唯一。Henniart 的策略大致如下：
\begin{enumerate}[label=(\roman*),leftmargin=2.2em]
\item 由 Bernstein--Zelevinsky（将在第 3 节讨论），可以归结为证明存在唯一一族双射
\[
\mathrm{rec}_n\colon\{\GL_n\ \text{超尖点表示的同构类}\}\longrightarrow
\{\text{$W_K$ 的}\ n\ \text{维不可约 WD 表示的同构类}\};
\]
\item 由 $\epsilon$-因子的性质（并同时考虑所有 $n$ 的 $\mathrm{rec}_n$），把上面的陈述
      归结为双射族 $\{\mathrm{rec}_n\}$ 的存在性；
\item 由 Henniart 的数值局部 Langlands 对应，映射 $\mathrm{rec}_n$ 的满性可由单射性推出；
\item 用 $L$-因子的性质证明这样的族 $\{\mathrm{rec}_n\}$ 必然单射。
\end{enumerate}
于是只需证明：在 $\GL_n$ 的不可约超尖点表示上存在满足定理中 (1)--(5) 的映射族
$\{\mathrm{rec}_n\}$。这正是 Harris--Taylor 证明、其后又由 Henniart 在 $K/\mathbb{Q}_p$
为局部数域时证明的陈述\en{译注}。看来 Scholze 的方法绕过了数值局部 Langlands 对应。
\end{remark}

\begin{yizhu}[底本的笔误：local number field]
底本此句作 “when $K/\mathbb{Q}_p$ is a local number field”，字面是「局部数域」，
但应与前一空格的 $\mathbb{Q}_p$ 一致，读作「$p$-进域」（$K/\mathbb{Q}_p$ 有限）。
译文保留「局部数域」并在方括号内标出正读。另：Harris--Taylor 与 Henniart 的原始结论
是对 $p$-进域上的 $\GL_n$ 的 LLC，这正是本文全篇的底域设定。
\end{yizhu}

\subsubsection{更精细的对应}

注意上面定理中包含的性质是为保证对应的唯一性所必需的（因而属于定理的一部分）。
下面几条可以看作额外的性质（它们对 $\mathrm{rec}_n$ 给出更精细的信息）：

\begin{theorem}
在上述局部 Langlands 对应（即定理中的 $\mathrm{rec}_n$）之下，有
\begin{enumerate}[label=(\arabic*),leftmargin=2.2em]
\item 不可约表示 $\pi\in\Repadm_G$ 超尖点，当且仅当 $\mathrm{rec}_n(\pi)$ 不可约；
\item $\pi$ 本质平方可积，当且仅当 $\mathrm{rec}_n(\pi)$ 不可分解；
\item $\pi$ 一般（generic），当且仅当 $L(\ad\circ\mathrm{rec}_n(\pi),s)$ 在 $s=1$ 处无极点。
\end{enumerate}
\end{theorem}

\subsubsection{惯性对应}

我们将在下一节陈述（$\GL_n(\mathcal{O}_K)$-）类型论的特例。虽然不能指望
$\GL_n(\mathcal{O}_K)$-表示与惯性表示之间存在对应，但两者之间有密切的联系。

\begin{proposition}
\leavevmode
\begin{enumerate}[label=(\arabic*),leftmargin=2.2em]
\item 设 $\pi$ 与 $\pi'$ 是 $G$ 的两个一般表示。若
\[
r_n(\pi)|_{I_K}\cong r_n(\pi')|_{I_K},
\]
其中 $r_n$ 记与 Weil--Deligne 表示 $\mathrm{rec}_n$ 相伴的 $W_K$ 的 $l$-进表示，则
\[
\pi|_{\GL_n(\mathcal{O}_K)}\cong\pi'|_{\GL_n(\mathcal{O}_K)}.
\]
\item （Paskunas）设 $\tau$ 是 $I_K$ 的 Weil--Deligne 惯性类型，且可延拓为 $W_K$ 的
      $n$ 维不可约 Frobenius 半单表示，则存在唯一（同构意义下）的光滑不可约表示
      $\sigma=\sigma_{\tau}$（$\GL_n(\mathcal{O}_K)$ 的），使得对 $\GL_n(K)$ 的任一光滑不可约
      无限维表示 $\pi$，
\[
\pi|_K\ \text{包含}\ \sigma_{\tau}
\quad\Longleftrightarrow\quad
\mathrm{rec}_n(\pi)|_{I_K}\cong\tau .
\]
\end{enumerate}
\end{proposition}

\begin{keytakeaways}
\begin{itemize}[leftmargin=1.5em]
\item LLC 就是一族双射 $\mathrm{rec}_n$：$\GL_n$ 的不可约容许表示 $\leftrightarrow$
      $W_K$ 的 $n$ 维 Frobenius 半单 Weil--Deligne 表示，由 (1)--(5) 五条性质唯一确定。
\item 唯一性由 Henniart 的论证保证：只要在超尖点表示上造出满足 (1)--(5) 的 $\mathrm{rec}_n$，
      其余由 Bernstein--Zelevinsky 分类自动拼装，且自动是双射。
\item 更精细的对应：超尖点 $\leftrightarrow$ 不可约；本质平方可积 $\leftrightarrow$ 不可分解；
      一般 $\leftrightarrow$ $L(\ad\circ\mathrm{rec},s)$ 在 $s=1$ 无极点。
\item 惯性对应说明 $\GL_n(\mathcal{O}_K)$ 上的表示（类型）与 $I_K$ 上的惯性类型之间有
      一一的联系，这是第 4 节的主题。
\end{itemize}
\end{keytakeaways}
